\documentclass[11pt]{article}
\usepackage{amsmath, amssymb, amsthm}
\usepackage{hyperref}

\theoremstyle{definition}
\newtheorem{definition}{Definition}[section]
\newtheorem{theorem}[definition]{Theorem}
\newtheorem{proposition}[definition]{Proposition}
\newtheorem{corollary}[definition]{Corollary}

\newcommand{\R}{\mathbb{R}}
\newcommand{\C}{\mathbb{C}}
\newcommand{\ReB}{\operatorname{Re}}
\newcommand{\ImB}{\operatorname{Im}}
\newcommand{\argB}{\operatorname{arg}}
\newcommand{\boxeq}{\boxed}
\newcommand{\rh}{\text{RH}}

\title{Spectral Fixed Point Theory 2.0: Core Derivation via a Zero-Free Wall}
\author{Ryan O. Van Gelder}
\date{December 30, 2025}

\begin{document}
\maketitle

\section{Preliminaries}
Let $\zeta(s)$ be the Riemann zeta function and let $\rho=\beta+i\gamma$ denote its non-trivial zeros (counted with multiplicity).
Define
\[
\beta_0 := \sup_{\rho}\ReB(\rho)\in \Big[\tfrac12,1\Big].
\]
The Riemann Hypothesis is equivalent to $\beta_0=\tfrac12$.

For $b\in[\tfrac12,1)$, we take $\argB\zeta(b+it)$ to mean the imaginary part of a branch of $\log\zeta$ defined along the vertical line $\Re(s)=b$ by continuous variation (with the standard convention that zeros/pole crossings are handled via the contour method used in \cite{SBM2009}).

\section{The SBM Master Identity and Instability Mass}

\subsection{Contour integral setup}
Sekatskii, Beltraminelli, and Merlini (SBM) \cite{SBM2009} consider a rectangular contour $C$ and the contour integral
\[
I=\oint_C g(z)\,\log\zeta(z)\,dz,
\qquad
g(z):=-i\frac{z-b}{\big(c^2-(z-b)^2\big)\big(d^2-(z-b)^2\big)},
\]
where $c,d>0$ and $c\neq d$. Evaluating $I$ by residues and comparing with its evaluation on the left edge $z=b+it$ yields an identity involving an integral of $\arg\zeta(b+it)$ and an explicit sum over zeros to the right of the wall $\Re(s)=b$.

\subsection{Identity with explicit zero term}
Define the weight function
\[
W_{c,d}(t) := \frac{t}{(c^2+t^2)(d^2+t^2)}.
\]
SBM obtain an equality of the form (their Eq.\,(1), specialized to $b<1$)
\begin{equation}\label{eq:SBM_master}
\int_{0}^{\infty}W_{c,d}(t)\,\argB\zeta(b+it)\,dt
=\mathcal{C}_{\mathrm{expl}}(b;c,d)-S(b;c,d),
\end{equation}
where $\mathcal{C}_{\mathrm{expl}}(b;c,d)$ is the explicit closed-form expression given below, and
\[
S(b;c,d)=\sum_{\rho:\ReB(\rho)>b}\Big(\text{explicit contribution from }\rho\Big),
\]
is a sum over \emph{exactly those zeros with $\ReB(\rho)>b$} (multiplicity included, conjugates paired). SBM emphasize that the sum is taken over zeros $\rho_k=\sigma_k+it_k$ with $b<\sigma_k$ and that conjugate pairs are combined to yield a real contribution. \cite[see Eq.\,(1) and surrounding discussion]{SBM2009}

\subsection{The explicit closed-form constant}
The constant $\mathcal{C}_{\mathrm{expl}}(b;c,d)$ is given by
\begin{equation}\label{eq:C_expl}
\boxeq{
\mathcal{C}_{\mathrm{expl}}(b;c,d)
:=
\frac{\pi}{2\,(d^2-c^2)}\,
\ln\!\left| \frac{\zeta(b+d)}{\zeta(b+c)} \right|
+
\frac{\pi}{2\,(d^2-c^2)}\,
\ln\!\left| \frac{\bigl(d^2-(1-b)^2\bigr)\,c^2}
                 {\bigl(c^2-(1-b)^2\bigr)\,d^2} \right|.
}
\end{equation}
This expression arises from the residues at the poles of the kernel $g(z)$ at $z=b+c$, $z=b+d$, and the contribution from the pole of $\zeta$ at $z=1$, as derived in \cite{SBM2009}.

\begin{remark}[On positivity/sign control]
SBM discuss sufficient conditions under which each off-wall zero contributes with the same sign, by comparing $c,d$ to the smallest ordinate $t_{\min}$ of any hypothetical zero with $\Re(\rho)>1/2$. They then give an explicit numerical admissibility regime for $(c,d)$ based on the verified height range available at the time of writing \cite{SBM2009}. This sign discussion is logically separate from the core equivalence (the integral equality holds iff there are no zeros with $\Re(\rho)>b$), but it motivates interpreting $S(b;c,d)$ as an ``instability mass.''
\end{remark}

\begin{definition}[Volchkov--Sekatskii divergence]
Fix admissible parameters $(c,d)$. Define the divergence as the absolute deviation from the SBM equality:
\[
\boxeq{
\Delta_V(b;c,d)
:=
\Bigl|\,
\mathcal{C}_{\mathrm{expl}}(b;c,d)
-
\int_{0}^{\infty}W_{c,d}(t)\,\argB\zeta(b+it)\,dt
\Bigr|.
}
\]
By \eqref{eq:SBM_master}, one has $\Delta_V(b;c,d)=|S(b;c,d)|$, and in the same-sign regime discussed by SBM one further has $\Delta_V(b;c,d)=S(b;c,d)\ge0$.
\end{definition}

\section{The theorem backbone: a zero-free wall}

\begin{theorem}[SBM zero-free wall equivalence \cite{SBM2009}]\label{thm:SBMwall}
Let $c,d>0$ with $c\neq d$, and let $b$ satisfy $1>b\ge\tfrac12$ and the admissibility hypotheses of SBM Theorem 1 (including the treatment of the limiting cases $b+c=1$, $b+d=1$). Then
\[
\boxeq{\Delta_V(b;c,d)=0\quad\iff\quad \nexists\,\rho\text{ with }\ReB(\rho)>b.}
\]
In particular,
\[
\boxeq{\RH\iff \Delta_V(\tfrac12;c,d)=0.}
\]
\end{theorem}

\section{SFPT equilibrium as a global, coercive wall}

\begin{definition}[SFPT spectral fixed point]
For fixed admissible $(c,d)$, define
\[
\mathcal{B}_{c,d}
:=
\{\,b\in[\tfrac12,1):\Delta_V(b;c,d)=0\,\},
\qquad
\boxeq{b_*:=\inf\bigl(\mathcal{B}_{c,d}\cup\{1\}\bigr).}
\]
\end{definition}

\begin{proposition}[Fixed point equals the supremum of zero real parts]
Let $\beta_0=\sup_{\rho}\ReB(\rho)$. For admissible $(c,d)$,
\[
\boxeq{b_*=\beta_0.}
\]
\end{proposition}

\begin{proof}
We consider two cases.
\begin{itemize}
    \item \textbf{Case $\beta_0=1$:} For every $b<1$, there exists a zero with $\ReB(\rho)>b$ (by definition of $\beta_0$). Hence, by Theorem~\ref{thm:SBMwall}, $\Delta_V(b;c,d)\neq 0$. Therefore $\mathcal{B}_{c,d}=\varnothing$ and $b_*=\inf\{1\}=1=\beta_0$.
    \item \textbf{Case $\beta_0<1$:} If $b>\beta_0$, then no zero satisfies $\ReB(\rho)>b$, so $\Delta_V(b;c,d)=0$ by Theorem~\ref{thm:SBMwall}. If $b<\beta_0$, then there exists a zero with $\ReB(\rho)>b$, so $\Delta_V(b;c,d)\neq 0$ by Theorem~\ref{thm:SBMwall}. Therefore $\mathcal{B}_{c,d}=[\beta_0,1)\cap[\tfrac12,1)=[\beta_0,1)$ and $b_*=\beta_0$.
\end{itemize}
\end{proof}

\begin{corollary}[RH as a global equilibrium]
\[
\boxeq{\RH\iff b_*=\tfrac12.}
\]
This is a statement of \textbf{global wall minimality} (``the first zero-free wall occurs at $1/2$''), not a symmetry-forced local stationarity condition.
\end{corollary}

\section{Computational SFPT core: prime--zero truncation mismatch}
The preceding sections are theorem-backed but ``zero-side.'' SFPT adds a computable proxy designed to converge to $\Delta_V$.

\subsection{Prime-side phase surrogate}
Define a smoothed truncation using von Mangoldt weights:
\[
\log \zeta_N(s) := \sum_{n \le N} \frac{\Lambda(n)}{n^s \log n}\, w_N(n),
\]
with a smoothing kernel $w_N$ (e.g.\ Ces\`aro or Gaussian). Define $\argB\zeta_N(b+it)$ by continuous phase tracking along $t$.

\subsection{Zero-side phase surrogate}
Given zeros $\rho$ with $|\ImB(\rho)|\le T$, define $\argB\zeta_T(b+it)$ via a truncated Hadamard/argument-principle representation (with the pole and $\Gamma$-factor handled consistently if one works with $\xi$ instead of $\zeta$).

\subsection{SFPT computable free-energy divergence}
\begin{definition}[SFPT free-energy divergence]
\[
\boxeq{
\Delta_{N,T}(b;c,d)
:=
\int_0^\infty W_{c,d}(t)\,
\bigl|\argB\zeta_N(b+it)-\argB\zeta_T(b+it)\bigr|^2\,dt.
}
\end{definition}

\subsection{SFPT convergence postulate (explicitly labeled)}
Under appropriate smoothing, branch conventions, and coupled limits $(N,T\to\infty)$,
\[
\Delta_{N,T}(b;c,d)\longrightarrow \Delta_V(b;c,d).
\]
This is the computational bridge: SFPT aims to approximate a known RH-equivalent obstruction by a measurable prime--zero mismatch, without assuming RH.

\section*{Proof sketch (contour $\to$ residues $\to$ zero-wall)}
The identity \eqref{eq:SBM_master} follows by applying the residue theorem to $I=\oint_C g(z)\log\zeta(z)\,dz$ with a rectangular contour $C$ whose left edge is $\Re(z)=b$ and whose right edge is taken to the right of $b+c$ and $b+d$ (and then sent to $\infty$ in the SBM argument). On the left edge $z=b+it$ one has $g(b+it)=\frac{t}{(c^2+t^2)(d^2+t^2)}\in\R$, so taking the imaginary part of the contour identity produces the weighted integral of $\arg\zeta(b+it)$. The poles of $g(z)$ at $z=b+c$ and $z=b+d$ produce the first logarithmic term in $\mathcal{C}_{\mathrm{expl}}(b;c,d)$, while the pole of $\zeta$ at $z=1$ contributes the second logarithmic term (handled via the generalized Littlewood theorem as in \cite{SBM2009}). Each nontrivial zero $\rho$ inside the contour contributes a term to $S(b;c,d)$, and in the limiting contour these are precisely the zeros with $\ReB(\rho)>b$. Thus $\Delta_V(b;c,d)=|S(b;c,d)|$ measures the total obstruction coming from zeros lying to the right of the wall $\Re(s)=b$.

\section{Numerical Methods}\label{sec:numerical}

This section specifies a concrete, reproducible protocol for computing the theorem-backed divergence
\(\Delta_V(b;c,d)\) and the SFPT proxy \(\Delta_{N,T}(b;c,d)\), together with grid rules, tolerances,
and diagnostics designed to detect branch, truncation-tail, and insufficient-spectrum artifacts.

\subsection{Canonical parameters and \(b\)-grid}\label{subsec:params}
Unless otherwise stated we fix
\[
(c,d)=\left(\tfrac32,\tfrac72\right),
\qquad
W_{c,d}(t)=\frac{t}{(c^2+t^2)(d^2+t^2)}.
\]
We evaluate on a coarse grid \(b\in\{0.55,0.60,0.65,0.70,0.75,0.80\}\) and then refine locally near any
observed transition or threshold in \(\Delta_V(b;c,d)\) or \(\Delta_{N,T}(b;c,d)\).
We avoid \(b=\tfrac12\) in initial numerics to prevent phase discontinuities at critical-line zeros.

\subsection{Computing \(\arg\zeta(b+it)\) without phase unwrapping}\label{subsec:arg}
To avoid branch/unwrap discontinuities, we compute \(\arg\zeta(b+it)\) from an ODE driven by the
logarithmic derivative. For \(s=b+it\),
\begin{equation}\label{eq:dargdt}
\frac{d}{dt}\arg\zeta(b+it)=\Re\!\left(\frac{\zeta'}{\zeta}(b+it)\right).
\end{equation}
We therefore define the continuous argument by
\begin{equation}\label{eq:arg_from_logderiv}
\arg\zeta(b+it)=\Arg(\zeta(b))+\int_{0}^{t}\Re\!\left(\frac{\zeta'}{\zeta}(b+i\tau)\right)d\tau,
\end{equation}
where \(\Arg(\zeta(b))\) is the principal argument of the real number \(\zeta(b)\) (i.e.\ \(0\) if
\(\zeta(b)>0\) and \(\pi\) if \(\zeta(b)<0\)); this pins the absolute phase at \(t=0\) and ensures
continuity in \(t\).

\subsection{\(t\)-grid and quadrature}\label{subsec:grid}
All \(t\)-integrals \(\int_0^\infty (\cdots)\,dt\) are truncated to \([0,T_{\max}]\) and evaluated with
composite Simpson's rule on a piecewise-uniform grid:
\[
\Delta t=
\begin{cases}
h_1=0.005, & t\in[0,5],\\
h_2=0.02, & t\in(5,50],\\
h_3=0.10, & t\in(50,T_{\max}].
\end{cases}
\]
Default \(T_{\max}=300\). Tail stability is assessed by repeating the computation with \(2T_{\max}\).
(Heuristically \(W_{c,d}(t)\sim t^{-3}\), so the weighted tail is typically mild, but we enforce this
numerically.)

\subsection{Gold-standard divergence \(\Delta_V(b;c,d)\)}\label{subsec:deltaV}
For each \(b\) we compute the explicit constant
\begin{equation}\label{eq:Cexpl_numeric}
\mathcal{C}_{\mathrm{expl}}(b;c,d)
=
\frac{\pi}{2(d^2-c^2)}
\left[
\ln\left|\frac{\zeta(b+d)}{\zeta(b+c)}\right|
+
\ln\left|\frac{(d^2-(1-b)^2)c^2}{(c^2-(1-b)^2)d^2}\right|
\right].
\end{equation}
We then compute \(\arg\zeta(b+it)\) on the grid via \eqref{eq:arg_from_logderiv} using
high-precision evaluation of \(\Re(\zeta'/\zeta)(b+it)\), and evaluate
\begin{equation}\label{eq:Ibcd}
I(b;c,d):=\int_{0}^{T_{\max}} W_{c,d}(t)\,\arg\zeta(b+it)\,dt.
\end{equation}
Finally,
\begin{equation}\label{eq:DeltaV_num}
\Delta_V(b;c,d):=\left|\mathcal{C}_{\mathrm{expl}}(b;c,d)-I(b;c,d)\right|.
\end{equation}

\paragraph{Recommended tolerances.}
We accept a computed \(\Delta_V(b;c,d)\) value only if all the following hold:
\begin{enumerate}
\item \textbf{Tail test:} \(|I_{2T_{\max}}-I_{T_{\max}}|\le \varepsilon_{\mathrm{tail}}\) with
\(\varepsilon_{\mathrm{tail}}=10^{-6}\) (tighten to \(10^{-8}\) for final plots).
\item \textbf{Mesh refinement:} halving each \(h_i\) changes \(I(b;c,d)\) by at most
\(\varepsilon_{\mathrm{mesh}}=10^{-6}\).
\item \textbf{Precision test:} increasing working precision (e.g.\ \(p\to p+40\) digits) changes
\(I(b;c,d)\) by at most \(\varepsilon_{\mathrm{prec}}=10^{-7}\).
\end{enumerate}

\subsection{SFPT proxy \(\Delta_{N,T}(b;c,d)\)}\label{subsec:deltaNT}
We define two continuous-phase surrogates \(\arg\zeta_N\) (prime side) and \(\arg\zeta_T\) (zero side),
both normalized at \(t=0\) to \(\Arg(\zeta(b))\).

\paragraph{Prime-side log-derivative surrogate.}
Fix a smoothing kernel; default is Gaussian
\[
w_N(n)=\exp\!\bigl(-(n/N)^2\bigr).
\]
Define
\begin{equation}\label{eq:logderiv_prime}
\left(\frac{\zeta'}{\zeta}\right)_N(s):=-\sum_{n\le N}\Lambda(n)\,n^{-s}\,w_N(n).
\end{equation}
Then set
\begin{equation}\label{eq:arg_prime}
\arg\zeta_N(b+it)
:=
\Arg(\zeta(b))
+
\int_0^t \Re\!\left(\frac{\zeta'}{\zeta}\right)_N(b+i\tau)\,d\tau.
\end{equation}

\paragraph{Zero-side log-derivative surrogate.}
Let \(\xi(s)=\tfrac12 s(s-1)\pi^{-s/2}\Gamma(s/2)\zeta(s)\).
Using a zero list \(\{\rho:|\Im(\rho)|\le T\}\), approximate
\begin{equation}\label{eq:logderiv_xi_trunc}
\left(\frac{\xi'}{\xi}\right)_T(s)\approx \sum_{|\Im(\rho)|\le T}\frac{1}{s-\rho}.
\end{equation}
Then define
\begin{equation}\label{eq:logderiv_zero}
\left(\frac{\zeta'}{\zeta}\right)_T(s)
:=
\left(\frac{\xi'}{\xi}\right)_T(s)
-\frac{1}{s}-\frac{1}{s-1}+\frac12\log\pi-\frac12\frac{\Gamma'}{\Gamma}\!\left(\frac{s}{2}\right),
\end{equation}
and set
\begin{equation}\label{eq:arg_zero}
\arg\zeta_T(b+it)
:=
\Arg(\zeta(b))
+
\int_0^t \Re\!\left(\frac{\zeta'}{\zeta}\right)_T(b+i\tau)\,d\tau.
\end{equation}

\paragraph{SFPT proxy divergence.}
With these surrogates we compute
\begin{equation}\label{eq:DeltaNT}
\Delta_{N,T}(b;c,d)
:=
\int_0^{T_{\max}} W_{c,d}(t)\,
\bigl|\arg\zeta_N(b+it)-\arg\zeta_T(b+it)\bigr|^2\,dt.
\end{equation}

\paragraph{Recommended schedules and tolerances.}
We recommend
\[
N\in\{10^4,10^5,10^6\},
\qquad
T\in\{T_{\max}+50,\;2T_{\max}+50\},
\]
with acceptance conditioned on:
\begin{enumerate}
\item \textbf{Normalization check:} \(\arg\zeta_N(b)=\arg\zeta_T(b)=\Arg(\zeta(b))\) (to numerical
roundoff).
\item \textbf{Monotone improvement:} for fixed \(b\), \(\Delta_{N,T}(b;c,d)\) should decrease (or at
least stabilize) as \(N\) and \(T\) increase past a transient regime.
\item \textbf{Tail test:} doubling \(T_{\max}\) changes \(\Delta_{N,T}\) by at most
\(\varepsilon^{(2)}_{\mathrm{tail}}=10^{-6}\).
\end{enumerate}

\subsection{Diagnostic checklist}\label{subsec:diagnostics}
Table~\ref{tab:diagnostics} lists common failure modes and how to resolve them.

\begin{table}[h]
\centering
\begin{tabular}{p{0.35\linewidth} p{0.28\linewidth} p{0.32\linewidth}}
\hline
\textbf{Symptom} & \textbf{Likely cause} & \textbf{Fix / test} \\
\hline
\(\arg\zeta(b+it)\) exhibits jumps of size \(\approx 2\pi\) &
Branch/unwrap bug (phase tracking) &
Use \eqref{eq:arg_from_logderiv} (integrate \(\Re(\zeta'/\zeta)\)) instead of direct \(\arg\) evaluation; verify continuity on a dense grid. \\
\hline
\(\Delta_V\) or \(I(b;c,d)\) changes materially when \(T_{\max}\) doubles &
Tail truncation bug &
Increase \(T_{\max}\); enforce \(|I_{2T_{\max}}-I_{T_{\max}}|\le \varepsilon_{\mathrm{tail}}\); consider extending the coarse-grid region with larger \(T_{\max}\). \\
\hline
\(\Delta_{N,T}\) does not improve as \(T\) increases &
Insufficient \(T\) (too few zeros / spectral truncation) &
Increase \(T\) (at least \(T\ge 2T_{\max}\)); spot-check \(\arg\zeta_T\) against a high-precision direct \(\arg\zeta\) at random \(t\)-points. \\
\hline
\(\Delta_{N,T}\) does not improve as \(N\) increases &
Insufficient \(N\) or poor smoothing &
Increase \(N\); use smoother \(w_N\) (Gaussian default); verify stability under modest changes of the smoothing scale. \\
\hline
Results change when increasing working precision &
Precision bug / cancellation &
Raise precision (e.g.\ \(p\to p+40\) digits) until \(\varepsilon_{\mathrm{prec}}\) is met; avoid subtracting nearly equal large quantities without extra precision. \\
\hline
Mesh halving changes outputs significantly &
Quadrature / grid resolution bug &
Halve \(\Delta t\) in each region; verify Simpson stability; densify near \(t\) where \(\Re(\zeta'/\zeta)\) is large. \\
\hline
\end{tabular}
\caption{Numerical diagnostics for the SBM-based SFPT 2.0 implementation.}
\label{tab:diagnostics}
\end{table}


\bibliographystyle{alpha}
\begin{thebibliography}{9}

\bibitem[SBM09]{SBM2009}
S.~K. Sekatskii, S.~Beltraminelli, D.~Merlini.
\newblock \emph{A few equalities involving integrals of the logarithm of the Riemann zeta-function and equivalent to the Riemann hypothesis II}.
\newblock arXiv:0904.1277 (2009).

\bibitem[Vol95]{Volchkov1995}
V.~V. Volchkov.
\newblock \emph{On an equality equivalent to the Riemann hypothesis}.
\newblock Ukrainian Mathematical Journal \textbf{47} (1995), 491--493.

\end{thebibliography}
\end{document}